\documentclass[a4paper,12pt,titlepage,finall]{article}

\usepackage[T1,T2A]{fontenc}
\usepackage[utf8]{inputenc}
\usepackage[russian]{babel}
\usepackage{tikz}
\usepackage{pgfplots}
\usepackage{geometry}
\usepackage{indentfirst}
\usepackage{listings}
\usepackage{amsmath}

\geometry{a4paper,left=30mm,top=30mm,bottom=30mm,right=30mm}
\setcounter{secnumdepth}{0}
\usepgfplotslibrary{fillbetween}
\pgfplotsset{compat=1.18}

\lstset{
    basicstyle=\ttfamily\small,
    breaklines=true,
    frame=single,
    inputencoding=utf8,
    extendedchars=true,
    keepspaces=true
}

\begin{document}

% Титульный лист
\begin{titlepage}
    \begin{center}
        {\small \sc Московский государственный университет \\
        имени М.~В.~Ломоносова \\
        Факультет вычислительной математики и кибернетики\\}
        \vfill
        {\Large \sc Отчет по заданию №6}\\
        ~\\
        {\large \bf <<Сборка многомодульных программ. \\
        Вычисление корней уравнений и определенных интегралов.>>}\\
        ~\\
        {\large \bf Вариант 3 / Ньютон+Хорды / Симпсон}
    \end{center}
    \begin{flushright}
        \vfill {Выполнил:\\
        студент 104 группы\\
        Коцин~А.~О.\\
        ~\\
        Преподаватель:\\
        % вставьте имя преподавателя
        }
    \end{flushright}
    \begin{center}
        \vfill
        {\small Москва\\2026}
    \end{center}
\end{titlepage}

\tableofcontents
\newpage

% ------------------------------------------------------------------
\section{Постановка задачи}

Требуется реализовать программу, вычисляющую с заданной точностью $\varepsilon$
площадь плоской фигуры, ограниченной тремя кривыми:
\[
    f_1(x) = e^{-x} + 3, \quad f_2(x) = 2x - 2, \quad f_3(x) = \frac{1}{x}.
\]

Площадь представляется как алгебраическая сумма определенных интегралов.
Пределы интегрирования --- абсциссы точек пересечения кривых, которые
находятся численно с точностью $\varepsilon_1$ двумя методами (выбор ---
через директиву препроцессора):
\begin{itemize}
    \item методом Ньютона (касательных),
    \item методом хорд (регула фальси).
\end{itemize}

Интегралы вычисляются с точностью $\varepsilon_2$ по формуле Симпсона
с правилом Рунге для контроля шага.

Функции $f_1$, $f_2$, $f_3$ и их производные реализованы на языке
ассемблера NASM (x86-32, соглашение \texttt{cdecl}).
Остальной код написан на языке Си. Сборка выполняется утилитой \texttt{make}.

Отрезки для поиска корней определяются аналитически: требуется, чтобы
$F(a) \cdot F(b) < 0$.

\newpage

% ------------------------------------------------------------------
\section{Математическое обоснование}

\subsection{Точки пересечения кривых}

\textbf{$f_2 \cap f_3$} --- единственный случай с точным решением:
\[
    2x - 2 = \frac{1}{x} \implies 2x^2 - 2x - 1 = 0
    \implies x_B = \frac{1 + \sqrt{3}}{2} \approx 1{,}3660.
\]

\textbf{$f_1 \cap f_3$} --- трансцендентное уравнение $e^{-x} + 3 - 1/x = 0$.
Анализ знаков:
\[
    F(0{,}20) = e^{-0{,}2} + 3 - 5{,}00 \approx -0{,}18 < 0, \quad
    F(0{,}35) = e^{-0{,}35} + 3 - 2{,}86 \approx +0{,}85 > 0.
\]
Корень $x_A \in [0{,}20;\; 0{,}35]$.

\textbf{$f_1 \cap f_2$} --- уравнение $e^{-x} + 5 - 2x = 0$.
Анализ знаков:
\[
    F(2{,}40) \approx +0{,}29 > 0, \quad F(2{,}70) \approx -0{,}33 < 0.
\]
Корень $x_C \in [2{,}40;\; 2{,}70]$.

\subsection{Условия сходимости методов}

Для методов Ньютона и хорд достаточно, чтобы $F'(x)$ не обращалась в нуль на отрезке~\cite{math}:

\begin{center}
\begin{tabular}{|l|l|c|}
\hline
Уравнение & $F'(x)$ & Знак на отрезке \\
\hline
$f_1 - f_3 = 0$ & $-e^{-x} + 1/x^2$ & $> 0$ на $[0{,}20;\;0{,}35]$ \\
$f_2 - f_3 = 0$ & $2 + 1/x^2$       & $> 0$ всегда \\
$f_1 - f_2 = 0$ & $-e^{-x} - 2$     & $< 0$ всегда \\
\hline
\end{tabular}
\end{center}

На каждом из трех отрезков $F$ строго монотонна, что гарантирует сходимость обоих методов.

\subsection{Выбор $\varepsilon_1$ и $\varepsilon_2$}

В точках пересечения подынтегральная функция $F(x) = f_i(x) - f_j(x)$
обращается в нуль. Поэтому погрешность $\delta x$ в определении границы интегрирования
вносит ошибку площади порядка $O(\delta x^2)$, что значительно меньше $\delta x$.
Таким образом, достаточно взять $\varepsilon_1 = \varepsilon_2 = \varepsilon$.

\subsection{Формула площади}

На интервале $[x_A, x_B]$: $f_1 > f_3 > f_2$, поэтому верхняя граница --- $f_1$, нижняя --- $f_3$.
На интервале $[x_B, x_C]$: $f_1 > f_2 > f_3$, поэтому верхняя --- $f_1$, нижняя --- $f_2$.
\[
    S = \int_{x_A}^{x_B} \bigl(f_1(x) - f_3(x)\bigr)\,dx
      + \int_{x_B}^{x_C} \bigl(f_1(x) - f_2(x)\bigr)\,dx.
\]

\subsection{Графики кривых}

\begin{figure}[h]
\centering
\begin{tikzpicture}
\begin{axis}[
    axis lines=middle,
    restrict x to domain=0.08:3.1,
    restrict y to domain=-0.5:5.5,
    enlargelimits,
    legend cell align=left,
    legend pos=north east,
    width=13cm, height=9cm,
    xlabel={$x$}, ylabel={$y$}]

\addplot[green!70!black, domain=0.08:3.1, samples=300, thick, name path=F1]
    {exp(-x)+3};
\addlegendentry{$f_1(x)=e^{-x}+3$}

\addplot[blue, domain=0.08:3.1, samples=300, thick, name path=F2]
    {2*x-2};
\addlegendentry{$f_2(x)=2x-2$}

\addplot[red, domain=0.08:3.1, samples=300, thick, name path=F3]
    {1/x};
\addlegendentry{$f_3(x)=1/x$}

\addplot[blue!15, samples=300]
    fill between[of=F1 and F3, soft clip={domain=0.2655:1.3660}];
\addplot[blue!15, samples=300]
    fill between[of=F1 and F2, soft clip={domain=1.3660:2.5395}];
\addlegendentry{$S \approx 3{,}636$}

\end{axis}
\end{tikzpicture}
\caption{Плоская фигура, ограниченная кривыми $f_1$, $f_2$, $f_3$}
\label{plot1}
\end{figure}

\newpage

% ------------------------------------------------------------------
\section{Результаты экспериментов}

Вычисленные координаты точек пересечения приведены в таблице~\ref{table1}.

\begin{table}[h]
\centering
\begin{tabular}{|c|c|c|}
\hline
Кривые & $x$ & $y$ \\
\hline
$f_1$ и $f_3$ & $0{,}2654743609$ & $3{,}7669$ \\
$f_2$ и $f_3$ & $1{,}3660254038$ & $0{,}7321$ \\
$f_1$ и $f_2$ & $2{,}5394547084$ & $3{,}0790$ \\
\hline
\end{tabular}
\caption{Координаты точек пересечения (метод Ньютона, $\varepsilon = 10^{-6}$)}
\label{table1}
\end{table}

Площадь фигуры при разных значениях точности (таблица~\ref{table2}):

\begin{table}[h]
\centering
\begin{tabular}{|c|c|c|c|}
\hline
$\varepsilon$ & Площадь & Итерации ($f_1{\cap}f_3$, $f_2{\cap}f_3$, $f_1{\cap}f_2$) & Метод \\
\hline
$10^{-4}$ & $3{,}6356614196$ & 3, 3, 2 & Ньютона \\
$10^{-6}$ & $3{,}6357850283$ & 3, 3, 3 & Ньютона \\
$10^{-8}$ & $3{,}6357851558$ & 4, 3, 3 & Ньютона \\
\hline
\end{tabular}
\caption{Результаты при различных значениях $\varepsilon$}
\label{table2}
\end{table}

\begin{figure}[h]
\centering
\begin{tikzpicture}
\begin{axis}[
    axis lines=middle,
    restrict x to domain=0.08:3.1,
    restrict y to domain=-0.5:5.5,
    enlargelimits,
    legend cell align=left,
    legend pos=north east,
    width=13cm, height=9cm,
    xticklabels={,,}, yticklabels={,,}]

\addplot[green!70!black, domain=0.08:3.1, samples=300, thick, name path=F1]
    {exp(-x)+3};
\addlegendentry{$f_1(x)=e^{-x}+3$}

\addplot[blue, domain=0.08:3.1, samples=300, thick, name path=F2]
    {2*x-2};
\addlegendentry{$f_2(x)=2x-2$}

\addplot[red, domain=0.08:3.1, samples=300, thick, name path=F3]
    {1/x};
\addlegendentry{$f_3(x)=1/x$}

\addplot[blue!20, samples=300]
    fill between[of=F1 and F3, soft clip={domain=0.2655:1.3660}];
\addplot[blue!20, samples=300]
    fill between[of=F1 and F2, soft clip={domain=1.3660:2.5395}];
\addlegendentry{$S \approx 3{,}6358$}

% Вертикальные пунктиры к точкам пересечения
\addplot[dashed, thin] coordinates { (0.2655, 3.7669) (0.2655, 0) };
\addplot[color=black] coordinates {(0.2655, 0)}
    node [label={below left:{\small $x_A$}}]{};

\addplot[dashed, thin] coordinates { (1.3660, 0.7321) (1.3660, 0) };
\addplot[color=black] coordinates {(1.3660, 0)}
    node [label={below:{\small $x_B$}}]{};

\addplot[dashed, thin] coordinates { (2.5395, 3.0790) (2.5395, 0) };
\addplot[color=black] coordinates {(2.5395, 0)}
    node [label={below right:{\small $x_C$}}]{};

\end{axis}
\end{tikzpicture}
\caption{Плоская фигура с отмеченными точками пересечения}
\label{plot2}
\end{figure}

\newpage

% ------------------------------------------------------------------
\section{Структура программы и спецификация функций}

Программа состоит из четырех модулей:

\begin{description}
    \item[\texttt{functions.asm}] ---
        функции $f_1$, $f_2$, $f_3$ и их производные $f_1'$, $f_2'$, $f_3'$,
        реализованные на ассемблере NASM x86-32 с использованием сопроцессора x87 FPU.
        Соглашение вызова --- \texttt{cdecl}: аргумент \texttt{long double}
        передается на стеке по смещению \texttt{[ebp+8]}, возвращаемое значение --- в \texttt{ST(0)}.

    \item[\texttt{root.c}] ---
        функция \texttt{root()}, находящая корень $F(x) = f(x) - g(x) = 0$ на отрезке $[a, b]$.
        Метод выбирается при компиляции:
        \texttt{-DMETHOD\_NEWTON} --- метод Ньютона,
        иначе --- метод хорд.

    \item[\texttt{integral.c}] ---
        функция \texttt{integral()}, вычисляющая определенный интеграл
        методом Симпсона с контролем точности по правилу Рунге.

    \item[\texttt{main.c}] ---
        разбор аргументов командной строки, вычисление корней и интеграла,
        вывод результатов, режим тестирования.
\end{description}

\medskip
\noindent Спецификации основных функций:

\begin{lstlisting}
/* Находит корень F(x) = f(x) - g(x) = 0 на [a, b] с точностью eps.
   df, dg - производные f и g (используются только методом Ньютона).
   *iters - счётчик итераций. */
long double root(
    long double (*f)(long double), long double (*g)(long double),
    long double (*df)(long double), long double (*dg)(long double),
    long double a, long double b, long double eps, int *iters);

/* Вычисляет интеграл f на [a, b] с точностью eps
   методом Симпсона (критерий Рунге). */
long double integral(
    long double (*f)(long double),
    long double a, long double b, long double eps);
\end{lstlisting}

\medskip
\noindent Схема зависимостей модулей:

\begin{center}
\begin{tikzpicture}[
    box/.style={draw, rounded corners, minimum width=3cm, minimum height=0.8cm,
                text centered, font=\ttfamily\small},
    arr/.style={->, thick}]
    \node[box] (asm)  at (0, 0)    {functions.asm};
    \node[box] (root) at (4, 0)    {root.c};
    \node[box] (intg) at (8, 0)    {integral.c};
    \node[box] (main) at (4, -2)   {main.c};
    \node[box] (prog) at (4, -4)   {prog};
    \draw[arr] (asm)  -- (prog);
    \draw[arr] (root) -- (prog);
    \draw[arr] (intg) -- (prog);
    \draw[arr] (main) -- (prog);
    \draw[arr] (main.north) -- (root.south);
    \draw[arr] (main.north) -- (intg.south);
\end{tikzpicture}
\end{center}

\newpage

% ------------------------------------------------------------------
\section{Сборка программы (Make-файл)}

Сборка выполняется командой \texttt{make} из директории \texttt{submissions/kotsin/}.
Для переключения метода нахождения корней используется параметр \texttt{METHOD}:

\begin{lstlisting}
make                       # метод Ньютона (по умолчанию)
make METHOD=METHOD_CHORDS  # метод хорд
make clean                 # удаление объектных файлов и исполняемого файла
\end{lstlisting}

\noindent Текст Make-файла:

\begin{lstlisting}
METHOD  ?= METHOD_NEWTON

CC       = gcc
ASM      = nasm
CFLAGS   = -m32 -Wall -Wextra -D$(METHOD)
ASMFLAGS = -f elf32

SRCDIR = src
OBJDIR = obj

OBJS = $(OBJDIR)/main.o $(OBJDIR)/root.o \
       $(OBJDIR)/integral.o $(OBJDIR)/functions.o

all: $(OBJDIR) prog

$(OBJDIR):
	mkdir -p $(OBJDIR)

prog: $(OBJS)
	$(CC) -m32 -no-pie -o $@ $^ -lm

$(OBJDIR)/main.o: $(SRCDIR)/main.c
	$(CC) $(CFLAGS) -I$(SRCDIR) -c -o $@ $<

$(OBJDIR)/root.o: $(SRCDIR)/root.c
	$(CC) $(CFLAGS) -I$(SRCDIR) -c -o $@ $<

$(OBJDIR)/integral.o: $(SRCDIR)/integral.c
	$(CC) $(CFLAGS) -I$(SRCDIR) -c -o $@ $<

$(OBJDIR)/functions.o: $(SRCDIR)/functions.asm
	$(ASM) $(ASMFLAGS) -o $@ $<

clean:
	rm -rf $(OBJDIR) prog
\end{lstlisting}

\newpage

% ------------------------------------------------------------------
\section{Отладка программы, тестирование функций}

Для запуска тестов используются ключи \texttt{--test-root} и \texttt{--test-integral}.
Выбор тестовых функций обеспечивает аналитически известный ответ и не более одного
тривиального теста (для Симпсона тривиальными являются полиномы степени $\leq 3$).

\subsection{Тестирование функции \texttt{root()}}

\textbf{Тест 1.} $x^2 = 2$ на $[1;\,2]$.

$F(x) = x^2 - 2$, $F'(x) = 2x > 0$ на $[1;\,2]$ (монотонно возрастает).
Точный корень: $x^* = \sqrt{2} \approx 1{,}41421356$.
$F(1) = -1 < 0$, $F(2) = 2 > 0$.

Результат: корень $= 1{,}414213562373095$, $|x - x^*| < 10^{-15}$, 4 итерации.

\medskip
\textbf{Тест 2.} $\sin x = 0$ на $[2;\,4]$.

$F(x) = \sin x$, $F'(x) = \cos x$.
На $[2;\,4]$: $F' < 0$ (монотонно убывает), единственный корень $x^* = \pi \approx 3{,}14159265$.
$F(2) > 0$, $F(4) < 0$.

Результат: корень $= 3{,}141592653589793$, $|x - x^*| < 10^{-15}$, 4 итерации.

\medskip
\textbf{Тест 3.} $e^x = 2$ на $[0{,}5;\,1]$.

$F(x) = e^x - 2$, $F'(x) = e^x > 0$ (монотонно возрастает).
Точный корень: $x^* = \ln 2 \approx 0{,}69314718$.
$F(0{,}5) \approx -0{,}35 < 0$, $F(1) \approx 0{,}72 > 0$.

Результат: корень $= 0{,}693147180559945$, $|x - x^*| < 10^{-15}$, 4 итерации.

\subsection{Тестирование функции \texttt{integral()}}

\textbf{Тест 1 (тривиальный).} $\displaystyle\int_0^1 x^2\,dx = \frac{1}{3}$.

Симпсон точен для полиномов степени $\leq 3$, поэтому ошибка машинная ($\approx 0$).
Результат: $0{,}333333333333333$, $|\Delta| < 10^{-15}$.

\medskip
\textbf{Тест 2.} $\displaystyle\int_0^{\pi} \sin x\,dx = 2$.

Результат: $2{,}000000000000984$, $|\Delta| \approx 10^{-12}$ при $\varepsilon = 10^{-12}$.

\medskip
\textbf{Тест 3.} $\displaystyle\int_1^{e} \frac{1}{x}\,dx = \ln e - \ln 1 = 1$.

Результат: $1{,}000000000000259$, $|\Delta| \approx 2{,}6 \cdot 10^{-13}$ при $\varepsilon = 10^{-12}$.

\newpage

% ------------------------------------------------------------------
\section{Программа на Си и на Ассемблере}

Исходные тексты программы приложены к отчету в архиве.
Архив содержит директорию \texttt{submissions/kotsin/} со следующей структурой:

\begin{lstlisting}
submissions/kotsin/
    Makefile
    src/
        functions.asm
        root.c      root.h
        integral.c  integral.h
        main.c
    report/
        report.pdf
        report.tex
\end{lstlisting}

\newpage

% ------------------------------------------------------------------
\section{Анализ допущенных ошибок}

Метод Ньютона сходится квадратично~\cite{math}: на гладких функциях с ненулевой
производной в корне каждая итерация удваивает число верных знаков.
Этим объясняется малое число итераций (3--4) даже при $\varepsilon = 10^{-8}$.

Метод хорд сходится сверхлинейно, что медленнее Ньютона: на тех же отрезках
требуется 3--11 итераций в зависимости от функции.

Для метода Симпсона погрешность убывает как $O(h^4)$, поэтому
правило Рунге с коэффициентом $1/15$ дает надежную оценку.
При $\varepsilon = 10^{-6}$ метод сходится за несколько удвоений шага.

Основной источник погрешности --- округление при вычислении $e^{-x}$
через цепочку инструкций x87 (\texttt{fldl2e}, \texttt{f2xm1}, \texttt{fscale}).
Погрешность этой последовательности не превышает нескольких ULP (единиц в
последнем разряде \texttt{long double}), что значительно лучше требуемой $\varepsilon = 10^{-8}$.

\newpage
\begin{raggedright}
\addcontentsline{toc}{section}{Список цитируемой литературы}
\begin{thebibliography}{99}
\bibitem{math} Ильин~В.\,А., Садовничий~В.\,А., Сендов~Бл.\,Х.
    Математический анализ. Т.\,1 ---
    Москва: Наука, 1985.
\end{thebibliography}
\end{raggedright}

\end{document}
